%==============================================================================
%  CAPÍTULO XIII
%==============================================================================
\chapter{Derecho concursal internacional e insolvencia transfronteriza}
\label{cap:transfronteriza}

\fichaCapitulo{Cooperación judicial y extraterritorialidad.}{Adopción de la Ley Modelo
UNCITRAL, determinación del \comi{} y análisis de los casos LATAM y Astaldi.}

%==============================================================================
\bloqueTeorico
%==============================================================================

\subsection{De la territorialidad al universalismo cooperativo}

El derecho concursal nació territorial. Cada Estado ejercía jurisdicción exclusiva sobre los
bienes situados en su territorio, y el concurso abierto en un país no producía efecto alguno
sobre los activos del deudor en otro. Este principio ---coherente con la soberanía--- resultaba
funcional cuando el patrimonio y la actividad del deudor coincidían con las fronteras.

La globalización de los mercados de capital volvió disfuncional ese esquema. Un deudor puede
hoy tener su administración en un país, sus activos productivos en otro y sus cuentas en un
tercero. Bajo la territorialidad estricta, el concurso se fragmenta en tantos procedimientos
como Estados, con administradores que compiten por los mismos bienes, acreedores que cobran
dos veces y deudores que aprovechan las costuras para ocultar activos.

El universalismo cooperativo ---el modelo que Chile adoptó--- no suprime la soberanía, pero la
articula: reconoce el procedimiento extranjero, coopera con el tribunal foráneo y coordina la
administración, de modo que el concurso opere como un todo aunque se desarrolle en varias
jurisdicciones.

\subsection{La Ley Modelo UNCITRAL en el Capítulo VIII de la Ley \Nro 20.720}

La adopción del modelo se produjo mediante la incorporación de la Ley Modelo de la CNUDMI
\parencite{ley-modelo-uncitral} como Capítulo VIII de la \LIR{}. La técnica ---trasladar un
texto uniforme internacional al derecho interno--- persigue previsibilidad: un representante
extranjero encuentra en Chile un procedimiento reconocible, análogo al de las demás
jurisdicciones que adoptaron el modelo.

El \art{302} confirma la vocación internacional del sistema: en caso de conflicto entre el
Capítulo VIII y una obligación nacida de un tratado en que Chile sea parte, \destacar{prevalece
el tratado}. Para un balance de la primera década de aplicación en Chile, véase
\textcite{doc-diez-anos-uncitral}.

%==============================================================================
\bloqueNormativo
%==============================================================================

\subsection{Ámbito de aplicación y autoridad competente}

El \art{300} delimita el ámbito del Capítulo VIII: se aplica cuando un tribunal o
representante extranjero solicita asistencia a las autoridades chilenas en relación con un
procedimiento extranjero, cuando se solicita asistencia en el exterior respecto de un
procedimiento chileno, o cuando concurren procedimientos paralelos.

El \art{303} distribuye la competencia: el reconocimiento lo ejercen los \destacar{tribunales
ordinarios de justicia}, los tribunales arbitrales cuando corresponda, y la \superir{} cuando
se hubiere iniciado una renegociación de la persona deudora. En materia de cooperación con
tribunales extranjeros intervienen, además, los administradores concursales cuando la
Superintendencia se lo requiera.

\subsection{Las definiciones del artículo 301: procedimiento principal y no principal}

El \art{301} define las categorías que estructuran todo el capítulo:

\begin{description}[leftmargin=0pt,style=nextline,font=\sffamily\bfseries]
  \item[Procedimiento extranjero] El procedimiento colectivo, judicial o administrativo,
  incluido el provisional, tramitado en un Estado extranjero con arreglo a una ley de
  insolvencia, en virtud del cual los bienes y negocios del deudor quedan sujetos al control o
  supervisión de un tribunal o representante extranjero, a efectos de reorganización o
  liquidación.

  \item[Procedimiento extranjero principal] El que se tramita en el Estado donde el deudor
  tiene el \destacar{centro de sus intereses principales} (\comi{}).

  \item[Procedimiento extranjero no principal] El que se tramita donde el deudor tiene un
  establecimiento, sin ser el centro de sus intereses principales.
\end{description}

\begin{foco}[Por qué la distinción principal / no principal lo decide todo]
El reconocimiento como procedimiento \emph{principal} produce efectos suspensivos automáticos y
potentes (\art{319}); el reconocimiento como \emph{no principal} sólo habilita medidas que el
tribunal debe conceder discrecionalmente. Determinar dónde está el \comi{} del deudor es, por
tanto, la cuestión central de todo litigio de reconocimiento, y el primer punto que debe
acreditar el representante extranjero.
\end{foco}

\subsection{Medidas provisionales desde la solicitud: el artículo 318}

El \art{318} permite al tribunal, \destacar{desde la presentación de la solicitud de
reconocimiento y hasta que se resuelva}, otorgar medidas provisionales cuando sean necesarias
y urgentes para proteger los bienes del deudor situados en Chile o los intereses de los
acreedores. Entre ellas, la suspensión de ejecuciones, la prohibición de transferir o gravar
activos locales y la entrega de información sobre el patrimonio.

\begin{practica}[La ventana provisional es la de mayor valor táctico]
Las medidas del \art{318} operan \emph{antes} del reconocimiento, cuando el activo local
todavía puede fugarse. Para el representante extranjero que persigue bienes ocultados en
Chile, solicitar estas medidas al presentar ---no después--- es lo que impide que el
patrimonio desaparezca durante la tramitación.
\end{practica}

\subsection{Efectos del reconocimiento del procedimiento principal: el artículo 319}

Reconocido un procedimiento extranjero \destacar{principal}, el \art{319} produce, mientras se
tramita, tres efectos automáticos:

\begin{enumerate}
  \item Se \textbf{suspende} el inicio o la continuación de las acciones y procedimientos
  individuales sobre los bienes, derechos, obligaciones o responsabilidades del deudor.
  \item Se \textbf{suspende} toda medida de ejecución contra sus bienes.
  \item Se \textbf{suspende} el derecho a transferir o gravar sus bienes o a disponer de
  ellos.
\end{enumerate}

Sobre la práctica chilena de reconocimiento, véase \textcite{doc-reconocimiento-extranjero};
para una guía de tramitación, \textcite{doc-transfronteriza-guia}.

\begin{alerta}[La suspensión del artículo 319 no equivale a la PFC del artículo 57]
Como se desarrolla en el bloque jurisprudencial, la doctrina administrativa ha precisado que
el reconocimiento de un procedimiento extranjero principal no habilita al juez chileno a
conceder automáticamente la \pfc{} general del \art{57} \Nro 1: los efectos se rigen por el
estatuto especial del Capítulo VIII, sin interpretaciones extensivas que perjudiquen a los
acreedores chilenos.
\end{alerta}

\subsection{El acceso directo del representante y de los acreedores extranjeros}

El modelo descansa sobre un principio de \destacar{acceso directo}. El \art{308} legitima al
representante extranjero reconocido para comparecer directamente ante los tribunales chilenos
---siempre por medio de abogado habilitado---, y el \art{309} precisa que la sola presentación de
la solicitud \destacar{no supone sumisión} del representante ni de los bienes del deudor en el
extranjero a la jurisdicción chilena para efecto distinto de la propia solicitud: una garantía
sin la cual ningún representante se arriesgaría a comparecer. Los \art{310} y \art{311} lo
facultan para pedir la apertura de un procedimiento local y para participar en el ya iniciado.

Por su parte, el \art{312} consagra la \destacar{igualdad de los acreedores extranjeros}: gozan
de los mismos derechos que los nacionales para iniciar el concurso y participar en él, sujetos al
mismo orden de prelación del Título XLI del Libro IV del \cc{}. La igualdad es procesal, no de
rango: el acreedor extranjero no salta la fila de las preferencias chilenas.

\subsection{El centro de los principales intereses (COMI) y su presunción}

La calificación del procedimiento como principal o no principal ---y con ella la intensidad de sus
efectos--- pende del \destacar{centro de los principales intereses} (COMI) del deudor. El
\art{315} \Nro 3 establece la presunción cardinal del sistema: \destacar{salvo prueba en
contrario, se presume que el domicilio social del deudor ---o su residencia habitual, si es
persona natural--- es el centro de sus principales intereses}. Es una presunción simplemente
legal: los acreedores locales pueden desvirtuarla acreditando que la administración efectiva del
deudor se sitúa en otro Estado, cuestión que en la práctica concentra el debate del
reconocimiento (véase el caso \textsc{latam} más adelante).

\subsection{Medidas posteriores al reconocimiento y protección de los acreedores}

Reconocido el procedimiento ---principal o no principal---, el \art{320} habilita un catálogo
\destacar{discrecional} de medidas más amplio que la suspensión automática: suspender acciones y
ejecuciones aún no paralizadas, examinar testigos y recabar información, encomendar al
representante la administración o realización de los bienes situados en Chile e, incluso,
\destacar{distribuir} esos bienes, siempre que el tribunal se asegure de que los acreedores
locales quedan suficientemente protegidos. Ese resguardo es explícito en el \art{321}, que impone
al tribunal velar por los intereses de acreedores y demás afectados ---incluido el deudor--- y lo
faculta para condicionar, modificar o revocar cualquier medida. La lógica es de contrapeso: a
mayor injerencia del procedimiento foráneo, mayor deber de tutela del interés local.

\subsection{Acciones revocatorias y cooperación entre tribunales}

El \art{322} legitima al representante extranjero, desde el reconocimiento, para \destacar{entablar
las acciones revocatorias concursales} de la ley chilena (capítulo \ref{cap:revocatorias}) ---con
la cautela, si el procedimiento es no principal, de que la acción recaiga sobre bienes que deban
administrarse en él---. Es la herramienta que impide que el deudor transfronterizo vacíe su
patrimonio local antes del reconocimiento.

El Título 4 erige el \destacar{deber de cooperación y comunicación directa}. Los \art{324} y
\art{325} obligan al tribunal y al administrador concursal chilenos a cooperar «en la medida de lo
posible» con los tribunales y representantes extranjeros, facultándolos para comunicarse
\destacar{directamente} ---sin exhorto internacional, según aclara el \art{314}---, con la sola
carga de publicar la comunicación en el \boletin{} dentro de dos días (cuya omisión no invalida lo
obrado). El \art{326} enumera las formas de cooperación: designación de personas coordinadoras,
intercambio de información, coordinación de la administración y aprobación de protocolos
concursales entre foros.

\subsection{Procedimientos paralelos y su coordinación}

Cuando coexisten un concurso local y uno extranjero sobre el mismo deudor, el sistema opta por la
\destacar{coordinación} antes que por la primacía. El \art{327} restringe la apertura de un
concurso chileno posterior al reconocimiento de uno extranjero principal: sólo procede si el deudor
tiene bienes en Chile, y sus efectos se \destacar{limitan a esos bienes} (concurso territorial
secundario). El \art{328} fija las reglas de compatibilidad ---la medida foránea cede ante el
concurso local en curso, y se reexamina si el local se abre después---, y el \art{329} ordena
coordinar varios procedimientos extranjeros entre sí. El principio rector, que el \art{307} manda
observar, es la interpretación uniforme y de buena fe propia de su \destacar{origen internacional}.

\begin{foco}[La cooperación no borra la preferencia local]
El Capítulo VIII abre las fronteras del concurso, pero no disuelve la tutela del acreedor chileno:
el COMI se presume en el domicilio pero admite prueba en contrario, las medidas discrecionales se
condicionan a la protección del interés local, y el concurso territorial del \art{327} reserva los
bienes situados en Chile. Universalismo, sí, pero \emph{cooperativo} y con resguardo del foro.
\end{foco}

%==============================================================================
\bloquePractico
%==============================================================================

\subsection{Solicitud de reconocimiento ante el juzgado civil chileno}

\begin{enumerate}
  \item \textbf{Presentación de la solicitud.} Escrito ante el juzgado de letras en lo
  civil competente, acompañando copia certificada de la resolución extranjera que declara
  la insolvencia y acredita la designación del representante extranjero.

  \item \textbf{Declaración de concursos paralelos.} Declaración jurada que individualice
  todos los procedimientos de insolvencia iniciados respecto del deudor de que tenga
  conocimiento el representante extranjero.

  \item \textbf{Medidas provisionales.} Solicitud incidental de paralización de las
  ejecuciones individuales dirigidas contra los activos locales y, en su caso,
  designación de un interventor.

  \item \textbf{Publicidad y oposición de terceros.} Publicación en el \boletin{} y
  sustanciación de las controversias de los acreedores locales.
\end{enumerate}

\subsection{Casos paradigmáticos de la práctica chilena}

La aplicación del Capítulo VIII, ante la ausencia de un reglamento detallado, ha generado
criterios interpretativos dispares. Conviene distinguir dos grandes grupos de casos: los de
\emph{persecución de activos chilenos en el extranjero} ---donde Chile es el foro principal---
y los de \emph{reconocimiento en Chile de procedimientos extranjeros}.

\begin{table}[htbp]
\centering
\footnotesize
\caption{Casos de insolvencia transfronteriza en la práctica chilena}
\label{tab:cap13-casos}
\begin{tabularx}{\textwidth}{@{}>{\raggedright\arraybackslash}p{0.15\textwidth}
                                 >{\raggedright\arraybackslash}p{0.24\textwidth} X@{}}
\toprule
\textbf{Caso} & \textbf{Tribunal / origen} & \textbf{Hitos y relevancia} \\
\midrule
Elimco Soluciones Integrales & 18.\textsuperscript{o} Juzgado Civil de Santiago; España
& Procedimiento principal en España (jul.\ 2015); solicitud en Chile el 26-05-2016;
\destacar{primer reconocimiento} del país el 02-06-2017, previo informe de la \superir{};
liquidación local de la agencia el 21-01-2020; cierre el 06-05-2022 \\
\addlinespace
Astaldi SpA & 11.\textsuperscript{o} y 19.\textsuperscript{o} Juzgado Civil de Santiago;
Italia & Juicios ejecutivos locales el 28-01-2019; en la audiencia inicial (25-02-2019) la
agencia opta por reorganización; acuerdo aprobado el 15-04-2019 y sancionado el 24-04-2019;
desglose a rol paralelo el 06-06-2019; modificaciones en 2020 y 2023 \\
\addlinespace
NBI Agencia en Chile & 22.\textsuperscript{o} Juzgado Civil de Santiago; Italia
& Solicitud de reconocimiento el 26-12-2018; informe de la \superir{} (31-01-2019) por no
precisar el objetivo; \destacar{archivo} el 11-11-2019. En paralelo, liquidación forzosa que
muta a reorganización simplificada sancionada el 26-05-2020 \\
\addlinespace
Isolux Ingeniería & 22.\textsuperscript{o} Juzgado Civil de Santiago, Rol C-3104-2018; España
& Reconocimiento el 16-11-2020, \destacar{sin juicio concursal local previo}; el tribunal
designó \emph{directamente} a los liquidadores locales, omitiendo el certificado de nominación
de la \superir{} (anomalía procesal) \\
\addlinespace
Assignia Infraestructuras & 10.\textsuperscript{o} Juzgado Civil de Santiago, Rol C-32054-2019;
España & Desistimiento del representante extranjero el 18-11-2021; oposición del liquidador
local (02-02-2022) por 972 acciones en el Consorcio de Salud Santiago Oriente y un arbitraje
pendiente; expone la \destacar{laguna} sobre la terminación del reconocimiento \\
\addlinespace
Ónix Capital / Chang Rajii & Rol C-22090-2016 (Grupo Arcano)
& El liquidador chileno obtuvo reconocimiento de la liquidación local en EE.UU., Reino Unido,
Isla de Man y Australia, repatriando activos por más de USD 10 millones \\
\addlinespace
LATAM Airlines Group & 2.\textsuperscript{o} Juzgado Civil de Santiago; EE.UU.\ (Cap.\ 11)
& Solicitud el 01-06-2020; oposición de acreedores por competencia y por el \art{301} (COMI en
Santiago); reconocimiento definitivo el 08-03-2023, con administrador local bajo la \superir{} \\
\bottomrule
\end{tabularx}
\end{table}

\verificar{Confirmar roles, fechas y estado procesal de cada caso contra los expedientes del
Poder Judicial antes de publicar.}

%==============================================================================
\bloqueJurisprudencial
%==============================================================================

\subsection{Lecciones dogmáticas de la práctica chilena}

El examen sistemático de estos expedientes revela la ausencia de un criterio de tramitación
homogéneo y tres problemas estructurales que el litigante debe conocer.

\subsubsection{La anomalía de Isolux: reconocimiento sin informe ni nominación}

En \falloc{Isolux Ingeniería}{C-3104-2018} el tribunal dictó el reconocimiento de la quiebra
española \destacar{designando directamente} a los liquidadores locales ---María Carolina Mira
Mora (titular) y Juan Pablo Lagos Rocafort (suplente)--- y decretando cautelares, sin solicitar
el informe de pertinencia de la \superir{} ni recabar el certificado de nominación. La anomalía
generó reparos institucionales sobre el alcance de las medidas de auxilio: el juez del
reconocimiento no puede sustituir el mecanismo legal de nominación de auxiliares.

\subsubsection{La laguna de Assignia: cómo se termina un reconocimiento}

\falloc{Assignia Infraestructuras}{C-32054-2019} expuso un vacío: ni la \LIR{} ni la Ley Modelo
regulan cómo \destacar{terminar o desistir} un procedimiento de reconocimiento ya otorgado. El
representante extranjero pretendió aplicar supletoriamente el \cpc{} para desistirse, alegando
que todos los activos locales estaban liquidados; el liquidador chileno se opuso, revelando que
el verdadero propósito era eludir la indisponibilidad que pesaba sobre 972 acciones del deudor
en el Consorcio de Salud Santiago Oriente y un arbitraje pendiente ante la Cámara de Comercio de
Santiago (condena de USD 500.000 a favor de Epsilon Consulting).

\begin{foco}[El reconocimiento no es una puerta que se cierra a voluntad]
Otorgado el reconocimiento y trabada la indisponibilidad sobre activos locales, el representante
extranjero no puede simplemente desistirse para liberarlos. Mientras subsistan bienes o juicios
locales, el liquidador chileno tiene interés y legitimación para oponerse. Es la salvaguarda de
los acreedores nacionales frente al vaciamiento por vía del desistimiento.
\end{foco}

\subsubsection{LATAM: el COMI y la oposición de los acreedores locales}

En \falloc{LATAM Airlines Group}{}\ la oposición de acreedores y accionistas minoritarios
evidenció la tensión entre extraterritorialidad y debido proceso local. Se alegó la incompetencia
de la justicia estadounidense para reestructurar los pasivos de una matriz chilena cuyo
\destacar{centro de intereses principales (COMI) estaba en Santiago}. El 2.\textsuperscript{o}
Juzgado Civil resolvió los incidentes de nulidad, confirmó la legitimación del Capítulo 11 y
ordenó designar un administrador local sujeto a las potestades de la \superir{}.

\subsection{Límites de las medidas de amparo locales}

De estos casos se decanta un principio que la doctrina administrativa ha confirmado: el
reconocimiento de un procedimiento extranjero principal \destacar{no faculta al juez chileno para
conceder automáticamente la \pfc{} general} del \art{57} \Nro 1. Los efectos se rigen por el
estatuto especial de medidas del Capítulo VIII (arts. 318-319), sin interpretaciones extensivas
que vulneren los privilegios de los acreedores chilenos disidentes.
